% =============================================================================
% SYSTEM INSTRUCTION TO LLM:
% This file is a PROMPT TEMPLATE.
% You MUST generate a FULL, FINAL, COMPILABLE LaTeX presentation
% based on this template.
%
% CRITICAL RULES:
% 1. DO NOT change anything in the Preamble.
% 2. You MUST delete all [STYLE REF] frames and the AI INSTRUCTION frame.
% 3. You MUST replace all placeholder texts with real content.
% 4. You MUST preserve the collegeBeamer theme and visual style.
% 5. Keep each slide's body text to roughly 8--10 lines, More than 10 pages is unacceptable.
% 6. NEVER use \includegraphics. No external image files exist. Draw every chart,
%    diagram and table with tikz / pgfplots / booktabs instead.
%    (The theme's own logo/background images are loaded by the .sty — not by you.)
% 7. Only these packages are available (already loaded in the preamble): collegebeamer,
%    tikz, pgfplots, amsmath, amssymb, pifont, booktabs, colortbl, multirow.
%    Do NOT add \usepackage lines.
% =============================================================================

\documentclass[aspectratio=169,10pt]{beamer}

% -----------------------------------------------------------------------------
% 1. Preamble & Theme (DO NOT CHANGE)
% -----------------------------------------------------------------------------
\usepackage[en,bjtu]{collegebeamer}

\usepackage{tikz}
\usetikzlibrary{positioning}
\usepackage{pgfplots}
\pgfplotsset{compat=1.17}
\usepackage{amsmath,amssymb}
\usepackage{pifont}
\usepackage{booktabs}
\usepackage{colortbl}
\usepackage{multirow}

% --- Font Setup ---
% Text: TeX Gyre Pagella (Palatino clone with real bold/italic under XeLaTeX,
% resolved from TeX Live by name). Math: newpxmath (Palatino-matched).
\usepackage{fontspec}
\setmainfont{texgyrepagella}[
    Extension = .otf,
    UprightFont = *-regular,
    BoldFont = *-bold,
    ItalicFont = *-italic,
    BoldItalicFont = *-bolditalic,
]
\let\mathbf\relax % avoid a redefinition clash between bm/beamer and newpxmath
\usepackage{newpxmath}

% CJK: route Chinese characters through serif/sans CJK faces so zh
% titles/authors/body never drop glyphs; Latin text keeps Pagella untouched.
% Prefer the Noto/Source Han faces (far better outlines) when installed --
% the Docker image ships fonts-noto-cjk -- and fall back to TeX Live's
% Fandol faces so a plain TeX Live still compiles.
\usepackage{xeCJK}
% Keep curly quotes Latin: xeCJK otherwise pads them to full width in English.
\xeCJKDeclareCharClass{Default}{"2018, "2019, "201C, "201D}
\IfFontExistsTF{Noto Serif CJK SC}{%
    \setCJKmainfont{Noto Serif CJK SC}[
        BoldFont = {* Bold},
        AutoFakeSlant = 0.15,
        SmallCapsFont = {Noto Serif CJK SC},
        SmallCapsFeatures = {},
    ]%
}{\IfFontExistsTF{Source Han Serif SC}{%
    \setCJKmainfont{Source Han Serif SC}[
        BoldFont = {* Bold},
        AutoFakeSlant = 0.15,
        SmallCapsFont = {Source Han Serif SC},
        SmallCapsFeatures = {},
    ]%
}{%
    \setCJKmainfont{FandolSong-Regular}[
        Extension = .otf,
        BoldFont = FandolSong-Bold.otf,
        AutoFakeSlant = 0.15,
        SmallCapsFont = FandolSong-Regular,
        SmallCapsFeatures = {},
    ]%
}}
\IfFontExistsTF{Noto Sans CJK SC}{%
    \setCJKsansfont{Noto Sans CJK SC}[
        BoldFont = {* Bold},
        AutoFakeSlant = 0.15,
    ]%
}{\IfFontExistsTF{Source Han Sans SC}{%
    \setCJKsansfont{Source Han Sans SC}[
        BoldFont = {* Bold},
        AutoFakeSlant = 0.15,
    ]%
}{%
    \setCJKsansfont{FandolHei-Regular}[
        Extension = .otf,
        BoldFont = FandolHei-Bold.otf,
        AutoFakeSlant = 0.15,
    ]%
}}

% Disable collegeBeamer auto TOC insertion
\AtBeginSection[]{}

% -----------------------------------------------------------------------------
% 1.1 Make "subtitle under each slide title" show FULL and stable
% -----------------------------------------------------------------------------
% Some collegeBeamer variants put only \insertsectionnumber under the title,
% which may appear as "0" or look incomplete.
% Here we redefine framesubtitle so it always shows:
%   "Section number  Section name"
% and avoids showing anything before the first \section is defined.
\setbeamertemplate{framesubtitle}{%
  \vspace{-0.6ex}%
  {\scriptsize\color{maincolor}%
    \ifnum\value{section}>0
      \thesection\quad\insertsectionhead
    \fi
  }\par
  \vspace{-0.4ex}%
}

% -----------------------------------------------------------------------------
% 2. Metadata (LLM MUST FILL)
% -----------------------------------------------------------------------------
% Keep title page compilable when subtitle/author/institute are left empty.
\setbeamertemplate{title page}{
    \hspace{-12mm}
    \begin{beamercolorbox}[wd=0.6\textwidth, sep=10pt, leftskip=8mm]{title}
        {\usebeamerfont{title}\inserttitle}\\
        \vskip5pt{\usebeamerfont{subtitle}\leavevmode\insertsubtitle}\\
        \vskip5pt{\usebeamerfont{author}\usebeamercolor[fg]{author}\leavevmode\insertauthor}\\
        \vskip2pt{\usebeamerfont{institute}\usebeamercolor[fg]{institute}\leavevmode\insertinstitute}\\
        \vskip5pt{\usebeamerfont{date}\usebeamercolor[fg]{date}\insertdate}
    \end{beamercolorbox}
}

\title[Short Title]{Full Presentation Title}
% Caller fills these three fields. Leave empty when unknown; never invent placeholders.
\subtitle{}
\author[]{}
\institute{}
\date{\today}

\begin{document}

% =============================================================================
% FRAME 1: TITLE PAGE
% =============================================================================
\maketitle

% =============================================================================
% FRAME 2: MANUAL TABLE OF CONTENTS
% =============================================================================
\begin{frame}{Outline}
\large
\renewcommand{\baselinestretch}{1.8}\selectfont

01 \quad Title of Part One

02 \quad Title of Part Two

03 \quad Title of Part Three

04 \quad Title of Part Four

05 \quad Summary \& Q\&A
\end{frame}

% =============================================================================
% SECTION 1 (EXAMPLE)
% IMPORTANT: define \section{...} before any content frames of that section
% so the subtitle line is complete (no "0").
% =============================================================================
\section{Title of Part One}

% =============================================================================
% FRAME 3: SECTION TRANSITION (EXAMPLE: SECTION 1)
% LLM MUST generate one such page before EACH section
% =============================================================================
\begin{frame}{Outline}
\large
\renewcommand{\baselinestretch}{1.8}\selectfont

\textbf{\textcolor{maincolor}{01 \quad Title of Part One}}

02 \quad Title of Part Two

03 \quad Title of Part Three

04 \quad Title of Part Four

05 \quad Summary \& Q\&A
\end{frame}

% =============================================================================
% CONTENT FRAME TEMPLATE (LLM MUST USE THIS STRUCTURE)
% =============================================================================
% Content page example 1: Text-left, figure-right layout
\begin{frame}{Content Example: Text Left, Figure Right}
\framesubtitle{} % keep empty; framesubtitle template will show section info automatically
\small
\begin{columns}[T]
\column{0.55\textwidth}
Write the core analysis content on the left side here.
This area is used for background description, problem definition,
logical reasoning, or methodology explanation.
Text should be concise and conclusion-oriented; avoid vague descriptions.

\vspace{0.6em}
\textbf{Key Takeaway:}
State the central message of this slide in one sentence.

\column{0.40\textwidth}
The right side is used for one of the following:
\begin{itemize}
\item Data table (booktabs + rowcolors)
\item TikZ flowchart
\item Checklist / comparison list
\end{itemize}
\end{columns}
\end{frame}

% Content page example 2: Text-top, figure-bottom layout
\begin{frame}{Content Example: Text Top, Figure Bottom}
\framesubtitle{} % automatically inherits section info
\small

% Upper text area
Write the high-level description or background introduction here.
This layout is ideal for displaying wide data charts, timelines,
system architecture overviews, or extended equation derivations.
Keep the text concise to leave room for the visual element below.

\vspace{1em}

% Lower figure/chart area
\centering
% \includegraphics[height=0.55\textheight, keepaspectratio]{example-image} % replace with your own image path
\begin{tikzpicture}
    % Placeholder diagram; replace with \includegraphics in practice
    \draw[fill=white, draw=maincolor, thick] (0,0) rectangle (10, 3.5);
    \node at (5, 1.75) {\textcolor{black}{Place wide chart / architecture diagram / timeline here}};
    \node[anchor=south east] at (10,0) {\tiny \textcolor{maincolor}{Source: Research Data 2024}};
\end{tikzpicture}

\vspace{0.5em}
\centerstate{
    \tiny \emph{Note: Supplementary remarks or data source attribution below the figure.}
}
\end{frame}


% Content page example 3: Figure-left, text-right layout
\begin{frame}{Content Example: Figure Left, Text Right}
\framesubtitle{} 
\small
\begin{columns}[c] % [c] = vertically centered; use [T] for top alignment if text is long

% Left: figure area (45\% width)
\column{0.45\textwidth}
    \centering
    % \includegraphics[width=\textwidth]{example-image} % replace in practice
    \begin{tikzpicture}
        \draw[fill=white, draw=maincolor, thick] (0,0) rectangle (5, 4);
        \node at (2.5, 2) {\textcolor{black}{Academic map / Lab equipment}};
    \end{tikzpicture}
    
% Right: text analysis area (50\% width)
\column{0.50\textwidth}
    \textbf{Figure / Object Analysis:}
    
    This section provides a detailed interpretation of the visual element
    on the left. Suitable for presenting:
    \begin{itemize}
        \item Physical structure of experimental apparatus
        \item Specific module details of an algorithm or model
        \item Domain-specific data analysis
    \end{itemize}

    \vspace{1em}
    \textbf{Key Findings:}
    \begin{enumerate}
        \item \ctextbf{Observation 1}: The figure on the left shows\ldots
        \item \ctextbf{Observation 2}: Compared to expectations, the results exhibit\ldots
    \end{enumerate}
\end{columns}
\end{frame}

% =============================================================================
% [STYLE REF] TABLE PAGE (DELETE IN FINAL OUTPUT)
% =============================================================================
\begin{frame}{[STYLE REF] Table Style Example}
\framesubtitle{}
\small
\centering
\renewcommand{\arraystretch}{1.15}
\footnotesize
\rowcolors{2}{maincolor!10}{white}
\begin{tabular}{p{0.2\textwidth} c c c}
\toprule
Metric & A & B & Conclusion \\
\midrule
Item 1 & 12 & 15 & Room for improvement \\
Item 2 & 20 & 18 & Stable \\
\bottomrule
\end{tabular}
\end{frame}

% =============================================================================
% [STYLE REF] TIKZ FLOW PAGE (DELETE IN FINAL OUTPUT)
% =============================================================================
\begin{frame}{[STYLE REF] Flowchart Style Example}
\framesubtitle{}
\centering
\begin{tikzpicture}[
process/.style={rectangle,draw=maincolor,fill=maincolor!10,
minimum height=1cm,text width=2cm,align=center}
]
\node[process] (a) {Step 1};
\node[process,right=of a] (b) {Step 2};
\node[process,right=of b] (c) {Step 3};
\draw[->,thick] (a)--(b);
\draw[->,thick] (b)--(c);
\end{tikzpicture}
\end{frame}

% =============================================================================
% [STYLE REF] TWO-COLUMN ANALYSIS (DELETE IN FINAL OUTPUT)
% =============================================================================
\begin{frame}{[STYLE REF] Two-Column Analysis Example}
\framesubtitle{}
\small
\begin{columns}
\column{0.52\textwidth}
Left column for argumentation, definitions, and derivations.

\column{0.43\textwidth}
Right column for action items, illustrations, or comparisons.
\end{columns}
\end{frame}

% =============================================================================
% AI INSTRUCTION PAGE (DELETE IN FINAL OUTPUT)
% =============================================================================
\begin{frame}[allowframebreaks]{AI GENERATION INSTRUCTIONS}
\framesubtitle{}
You MUST:
\begin{itemize}
\item Delete all [STYLE REF] frames and this instruction frame
\item Replace all placeholders with real content
\item Insert a highlighted TOC transition page before each section
\item Define \texttt{\textbackslash section\{...\}} before the first frame of that section
\item Use columns layout for all content pages
\item Keep information density high (no empty slides)
\item Output FULL LaTeX code only
\end{itemize}
\end{frame}

% =============================================================================
% FINAL PAGE
% =============================================================================
\QApage

\end{document}
